% Generated draft: requires mathematical and source review.
\documentclass[11pt]{article}
\usepackage[margin=25mm]{geometry}
\usepackage{fontspec}
\setmainfont{Latin Modern Roman}
\usepackage{amsmath,amssymb,mathtools,unicode-math}
\setmathfont{Latin Modern Math}
\usepackage{enumitem,needspace,xurl,hyperref}
\setlength{\parindent}{0pt}
\setlength{\parskip}{5pt}
\setlength{\emergencystretch}{3em}
\newcommand{\N}{\mathbb{N}}
\newcommand{\Z}{\mathbb{Z}}
\newcommand{\Q}{\mathbb{Q}}
\newcommand{\R}{\mathbb{R}}
\newcommand{\C}{\mathbb{C}}
\newcommand{\F}{\mathbb{F}}
\newcommand{\A}{\mathbb{A}}
\newcommand{\PP}{\mathbb{P}}
\newcommand{\E}{\mathbb{E}}
\newcommand{\eps}{\varepsilon}
\newcommand{\dd}{\,\mathrm{d}}
\begin{document}
\section*{UnsolvedMath ID: 30001356}
\textbf{Source problem code:} catalogue entry
\subsection{Definitions and mathematical statement}
A smooth embedded surface \(M\subset\R^3\) is minimal if its mean curvature vanishes identically, and complete if its induced Riemannian metric is complete. Non-planar means not contained in an affine plane. Ends are topological ends, described by the unbounded components remaining outside an exhaustion by compact subsets. Under the connected, orientable, boundaryless finite-topology convention, a surface of genus \(g\) with \(n\) ends is homeomorphic to \(\Sigma_g\setminus\{p_1,\ldots,p_n\}\), where \(\Sigma_g\) is a compact connected orientable surface of genus \(g\) and the punctures are distinct.

The supplied record asks two questions. Its uniqueness relation and several surface conventions are unspecified, so the following is a conditional formulation rather than a fully determined conjecture. Connectedness, orientability, and absence of boundary are used explicitly as provisional conventions; no properness hypothesis is added.

\begin{enumerate}
\item For each \(g\in\Z_{\geq0}\), let \(\mathcal C_g\) denote the class of connected, orientable, boundaryless, non-planar complete embedded minimal surfaces in \(\R^3\) of genus \(g\) with exactly one end. Does
\[
\forall g\in\Z_{\geq0},\qquad
\mathcal C_g\ne\varnothing
\quad\text{and}\quad
\forall M,N\in\mathcal C_g,\ M\sim N
\]
hold? Here \(\sim\) must be an equivalence relation expressing the intended meaning of uniqueness. The source does not specify whether it allows ambient isometries, dilations, or another equivalence, so this part remains a statement schema until that choice is supplied. The first question does not explicitly impose finite topology.

\item Under the finite-topology convention above, for every \(g\in\Z_{\geq0}\) and every integer \(n\geq2\), is the following equivalence true?
\[
\begin{gathered}
\text{There exists a smooth embedding }
 f:\Sigma_g\setminus\{p_1,\ldots,p_n\}\longrightarrow\R^3\\
\text{whose induced metric is complete and whose mean curvature vanishes}
\end{gathered}
\quad\Longleftrightarrow\quad n\leq g+2.
\]
This interprets the second question as existence for a topological type. No conformal structure on the punctured surface is supplied or prescribed. The assertion concerns both necessity and sufficiency of the inequality.
\end{enumerate}

The source labels these questions open. This is a supplied status label, not an independently verified mathematical conclusion.
\subsection{Short English statement}
The record asks whether every nonnegative genus has a unique non-planar complete embedded minimal surface in \(\R^3\) with one end, and whether a finite-topology surface of genus \(g\) with \(n\geq2\) ends admits a complete minimal embedding exactly when \(n\leq g+2\). Uniqueness requires an unspecified equivalence relation, and the surface conventions require confirmation.
\subsection{Sources}
\begin{itemize}
\item Local UnsolvedMath registry: supplied complete record, permanent ID \(30001356\), title \emph{Finite-Topology Embedded Minimal Surfaces}, exact source problem code \texttt{catalogue entry}, collection \emph{Oberwolfach Reports Open Problems}. No individual proposer is supplied; the proposed year is recorded as \(2009\).
\item Supplied report citation: \emph{Arbeitsgemeinschaft: Minimal Surfaces (2010)}, Oberwolfach Reports, pp.~2545--2584, DOI \texttt{10.4171/owr/2009/45}. \url{https://doi.org/10.4171/owr/2009/45}. The record describes Oberwolfach Reports as documenting mathematical workshops and their open-problem discussions. The supplied title contains \(2010\), whereas the proposed year and DOI contain \(2009\); this has not been resolved.
\item Supplied literature source, labelled ``OWR report, stated open conjectures'': \url{https://ems.press/journals/owr/articles/4128}.
\item Supplied literature source, labelled ``Weber--Hoffman--Wolf, genus-one helicoid'': \url{https://annals.math.princeton.edu/2009/169-2/p01}.
\item Supplied literature source, labelled ``original OWR report'': \url{https://doi.org/10.4171/owr/2009/45}.
\item Literature assessment checked \texttt{2026-08-21}: supplied dated curation, not independently verified. It labels the problem open and states that uniqueness for one-ended surfaces of each genus and the Hoffman--Meeks end bound remain open in general, while genus-one existence is known and does not settle uniqueness or existence for all genera. The background calls this dated literature triage derived from the supplied OWR curation and requests independent verification of time-sensitive status and cited claims.
\item The registry marks the statement \texttt{corrected\_verified}, but its accompanying verification only reports adding mathematical delimiters without changing the mathematical content. This does not establish mathematical verification of the formulation.
\end{itemize}
\end{document}
